\documentclass[11pt]{article}
\usepackage[a4paper,margin=1in]{geometry}
\usepackage{amsmath,amssymb,amsthm,booktabs,hyperref}
\hypersetup{colorlinks=true,linkcolor=blue,urlcolor=blue}
\newtheorem{theorem}{Theorem}
\newcommand{\E}{\mathcal E}
\title{The maximum-run enumeration of flat Fubini rankings}
\author{Principia Math in Progress}
\date{12 September 2026}
\begin{document}
\maketitle

\paragraph{Task 1: Resolve MathDB 367940.}
Let \(s(n,r)\) count Fubini rankings of length \(n\) whose maximal
strictly increasing runs have strictly increasing first entries,
with exactly \(r\) runs. We prove the maximum-run enumeration
proposed in Section 4.1 of Barrese, Elder, Harris, and Simpson~\cite{source}.
Their displayed equation has indexing errors, which we make explicit
before proving the intended enumeration for every size.

Use the second-order Eulerian convention of the cited OEIS entry
A340556~\cite{eulerian}:
\begin{equation}\label{erec}
 E_2(0,0)=1,\qquad
 E_2(j,k)=kE_2(j-1,k)+(2j-k)E_2(j-1,k-1),
\end{equation}
with entries outside \(0\le k\le j\) equal to zero.
Write \(\E_j(z)=\sum_{k=0}^jE_2(j,k)z^k\).
For \(j\ge1\), this counts descents of Stirling permutations
\emph{including the final descent to an appended zero}.

\begin{theorem}\label{main}
The largest possible number of runs in a flattened Fubini ranking
of length \(2j+1\) is \(j+1\). For every \(j\ge0\),
\begin{equation}\label{answer}
 \boxed{\quad s(2j+1,j+1)=\E_j(2)
             =\sum_{k=0}^j E_2(j,k)2^k.\quad}
\end{equation}
More precisely, the number of these rankings with exactly \(k+1\)
distinct ranks is
\begin{equation}\label{refinement}
 N_j(k)=\sum_{m=0}^k(-1)^{k-m}\binom{k}{m}e_j(1,2,\ldots,m),
 \qquad 0\le k\le2j,
\end{equation}
where \(e_j\) is the elementary symmetric polynomial, with
\(e_0=1\) and \(e_j(1,\ldots,m)=0\) for \(j>m\).
\end{theorem}

\paragraph{The indexing issue.}
The source's Conjecture 4.4, equation (4.1) in the publisher's current
accepted manuscript, prints
\[
 s(2j+1,j-1)=\sum_{i=0}^{j-1}E_2(2j+1,i)2^i.
\]
With its cited A340556 convention this literal formula is false:
at \(j=2\) its left side is \(s(5,1)=1\), whereas its right side
is \(2\). The surrounding text instead specifies odd length and
the \emph{largest possible} number of runs, and its Table 2 gives
\[
 s(3,2)=2,\quad s(5,3)=10,\quad s(7,4)=82,\quad s(9,5)=938.
\]
It identifies this diagonal with A112487, whose defining formula is
\(\E_j(2)\)~\cite{sequence}. Thus~\eqref{answer} is precisely the
enumeration specified by that text, table, and sequence. The proof
below resolves that enumeration; noticing the printed mismatch
alone is not the result.

\paragraph{1. Replace ranks by their order.}
A \emph{packed word} uses exactly an alphabet \([q]=\{1,\ldots,q\}\).
There is a bijection from Fubini rankings to packed words: replace
the smallest occurring rank by \(1\), the next by \(2\), and so on.
Conversely, if letter \(i\) has multiplicity \(c_i\), replace every
\(i\) by
\[
 1+c_1+\cdots+c_{i-1}.
\]
These are exactly the valid competition ranks with ties.
Both maps preserve strict comparisons and equalities, hence
preserve all the run conditions.

No run except the last can have length one: its single entry
would be smaller than the next run's first entry, so the two
would belong to the same increasing run. Consequently a word
with \(r\) runs has length at least \(2r-1\).
At length \(2j+1\), equality with \(r=j+1\) forces the run lengths
to be \(2,2,\ldots,2,1\). Such a packed word has the form
\begin{equation}\label{word}
 a_0\,b_1\mid a_1\,b_2\mid\cdots\mid a_{j-1}\,b_j\mid a_j,
 \qquad
 a_0<a_1<\cdots<a_j,\quad b_i\ge a_i.
\end{equation}
Conversely these inequalities make the displayed bars exactly
the run boundaries. They also force \(a_0\) to be the unique
smallest letter.

Remove \(a_0\) and decrease all remaining packed labels by one.
The resulting objects are precisely the surjective assignments
to \(2j\) named positions satisfying
\begin{equation}\label{poset}
 1\le a_1<\cdots<a_j\le k,\qquad a_i\le b_i\le k.
\end{equation}
Any such assignment extends uniquely back to~\eqref{word} by
increasing all labels and adjoining the new minimum.
For example, choosing \(a_i=b_i=i\) gives a valid object for
every \(j\), so the run bound is attained.
The case \(j=0\) is the unique one-letter ranking.

\paragraph{2. Count assignments before requiring surjectivity.}
Let \(T_j(m)\) count assignments satisfying~\eqref{poset} into
\([m]\), now allowing unused labels. After choosing the increasing
sequence of \(a_i\), each \(b_i\) has \(m-a_i+1\) choices. Therefore
\begin{equation}\label{elementary}
 T_j(m)=\sum_{1\le a_1<\cdots<a_j\le m}
                 \prod_{i=1}^j(m-a_i+1)
        =e_j(1,2,\ldots,m).
\end{equation}
Choosing the used \(k\)-element subalphabet and packing it gives
\begin{equation}\label{binomial}
 T_j(m)=\sum_{k=0}^{2j}\binom{m}{k}N_j(k).
\end{equation}
Binomial inversion proves~\eqref{refinement}.
Moreover, the identity
\[
 \sum_{m=k}^{\infty}\frac{\binom{m}{k}}{2^{m+1}}=1
\]
gives the useful averaging formula
\begin{equation}\label{average}
 s(2j+1,j+1)=\sum_{k=0}^{2j}N_j(k)
            =\frac12\sum_{m=0}^{\infty}
                    e_j(1,2,\ldots,m)\,2^{-m}.
\end{equation}
There is no convergence issue: for fixed \(j\) the symmetric
polynomial is bounded by a constant times \(m^{2j}\).

\paragraph{3. Evaluate the average.}
We give the generating-function calculation explicitly.
Put
\[
 H_j(t)=\sum_{m\ge0}e_j(1,2,\ldots,m)t^m,\qquad |t|<1.
\]
Adjoining the largest argument \(m\) in an elementary symmetric
polynomial gives
\[
 e_j(1,\ldots,m)=e_j(1,\ldots,m-1)
                       +m e_{j-1}(1,\ldots,m-1).
\]
Thus \(H_0(t)=(1-t)^{-1}\) and, for \(j\ge1\),
\begin{equation}\label{hrec}
 H_j(t)=\frac{t}{1-t}\frac{d}{dt}\bigl(tH_{j-1}(t)\bigr).
\end{equation}
If \(H_j(t)=t^jQ_j(t)/(1-t)^{2j+1}\), this recurrence becomes
\begin{equation}\label{qrec}
 Q_0(t)=1,\qquad
 Q_j(t)=\bigl(j+(j-1)t\bigr)Q_{j-1}(t)
                      +t(1-t)Q_{j-1}'(t).
\end{equation}

For completeness, the Eulerian recurrence~\eqref{erec} follows by
inserting the adjacent largest pair \(jj\) in a Stirling permutation
of order \(j-1\). There are \(2j-1\) insertion gaps. If the old
permutation has \(h\) descents, including the terminal descent,
exactly \(h\) gaps preserve the count, and the other gaps increase
it by one. This also covers \(j=1\), starting from the empty word.
In polynomial form the recurrence is
\[
 \E_j(z)=(2j-1)z\E_{j-1}(z)
                         +z(1-z)\E_{j-1}'(z).
\]
Substitution shows that \(t^j\E_j(1/t)\) satisfies~\eqref{qrec}
with the same initial value. Hence
\[
 Q_j(t)=t^j\E_j(1/t),\qquad
 H_j(t)=\frac{t^{2j}\E_j(1/t)}{(1-t)^{2j+1}}.
\]
Finally~\eqref{average} gives
\[
 s(2j+1,j+1)=\frac12H_j(1/2)=\E_j(2),
\]
as asserted. This proves the formula and the refinement for
every \(j\ge0\).
\hfill\(\square\)

\paragraph{Verification and scope.}
The exact verifier passes 120 checks, including independent enumeration
of all packed words through length \(7\), the conversion back to valid
Fubini rankings, the refined counts, and rational generating-function
identities. The first diagonal values are
\[
 1,\ 2,\ 10,\ 82,\ 938,\ 13778,\ 247210,\ 5240338,\ 128149802.
\]
The proof, rather than these finite checks, establishes the infinite
identity. The general two-variable enumeration and the source's
separate total-count conjecture are not asserted here.
The dated status audit located no matching full public claim and
confirmed that the current MathDB record remains open.

\paragraph{Task 2: Construct a direct bijection to decorated Stirling permutations.}
Replace the averaging argument by a bijection between maximum-run
rankings and Stirling permutations in which each descent, including
the terminal descent, receives one of two colors.

\paragraph{Task 3: Enumerate the maximum-run rankings of even length.}
Determine \(s(2j,j)\), keeping track of the possible run-length
profiles and the number of distinct ranks.

\paragraph{Task 4: Enumerate all run counts.}
Extend the packed-word description to derive a recurrence or generating
function for \(s(n,r)\) for arbitrary \(n,r\), and compare its total
count with the source's separate generating-function conjecture.

\paragraph{Task 5: Determine the distribution of distinct ranks.}
Choose uniformly from the maximum-run rankings of length \(2j+1\).
Use~\eqref{refinement} to find the asymptotic mean, variance, and
limiting distribution for the number of distinct ranks.

\begin{thebibliography}{9}
\bibitem{source}
K. Barrese, J. Elder, P. E. Harris, and A. Simpson,
\emph{Enumerating flat Fubini rankings},
The Art of Discrete and Applied Mathematics, accepted manuscript
published 25 May 2026, DOI 10.26493/2590-9770.1883.9e1.
Conjecture 4.4 and Table 2, page 19.
\url{https://adam-journal.eu/index.php/ADAM/article/view/1883}.
Author preprint: \url{https://arxiv.org/abs/2504.06466}.
\bibitem{eulerian}
The OEIS Foundation, A340556, second-order Eulerian numbers,
including the terminal descent.
\url{https://oeis.org/A340556}.
\bibitem{sequence}
The OEIS Foundation, A112487, values \(\E_j(2)\).
\url{https://oeis.org/A112487}.
\end{thebibliography}
\end{document}
